\documentclass[10pt,a4paper]{amsart}
\usepackage[margin=19mm]{geometry}
\usepackage{amsmath,amssymb,mathtools}
\usepackage[scr=rsfs,cal=euler]{mathalfa}
\usepackage{xfrac}
\usepackage[unicode,hidelinks,bookmarks=false]{hyperref}
\newcommand{\ie}{i.e.\ }
\newcommand{\cf}{cf.\ }
\newcommand{\resp}{respectively\ }
\newcommand{\Subsection}{\subsection}
\providecommand{\nobreakdash}{\nobreak\hbox{-}}
\setlength{\emergencystretch}{2em}
\multlinegap0pt
\usepackage{bm}
\usepackage{bbm}
\usepackage{dsfont}
\usepackage{xspace}
\usepackage{cancel}
\usepackage{mathtools}
\usepackage{amsthm}
\makeatletter
\@ifundefined{thm}{\newtheorem{thm}{Theorem}}{}
\@ifundefined{prop}{\newtheorem{prop}[thm]{Proposition}}{}
\makeatother
\title[Reduction of Poisson manifolds with Hamiltonian Lie algebroi]{Reduction of Poisson manifolds with Hamiltonian Lie algebroids}
\author{Yuji Hirota, Noriaki Ikeda}
\date{Source: arXiv:2411.10969v3 (CC BY 4.0)}
\begin{document}
\maketitle

\noindent\emph{Source and licence.} Reduction of Poisson manifolds with Hamiltonian Lie algebroids. Version arXiv:2411.10969v3 (CC BY 4.0), distributed under the Creative Commons Attribution 4.0 International licence, as recorded in the arXiv licence metadata for this exact version and shown on the abstract page https://arxiv.org/abs/2411.10969v3. Full rights notice: licenses/arxiv-2411.10969-CC-BY-4.0.txt.\par\smallskip

\noindent\textit{Editorial scope.} Counted unit: the Thm whose source label is \texttt{sec3:main thm} in the article (source0.tex lines 348--363, source label \texttt{sec3:main thm}), delivered together with the article's own complete original proof of it (source0.tex lines 364--417, one \texttt{proof} environment of 54 lines). No mathematics is written, repaired or re-derived: statement and proof are the article's own text line for line. Required context retained uncounted: sec3:prop\_vanishing (lines 297--300). Every definition, display and label of those lines is retained unchanged. Omitted from this package and not redistributed: 1--296, 301--347, 418--550. Nothing inside a retained excerpt is omitted or altered: every retained body is delivered exactly as the article's lines are written, so no omission receipt of any kind accompanies this package. Citations: the retained excerpts carry no citation command at all, so no bibliography entry of the article is transcribed and no reference is redistributed. Reference adaptation: none was needed and none was made; every reference inside the delivered unit resolves inside it, so no unresolved reference or citation is produced. Printed slips kept literal: no printed word, symbol, hypothesis or number of the counted statement or of its proof was altered. Layout and class: the article's own class is \texttt{article} with packages including geometry, amsmath, amssymb, amsfonts, amscd, amsthm, ascmac, eucal, enumerate, indentfirst, graphicx, color; the wrapper is the lane's pinned \texttt{amsart} editorial wrapper at 10pt on A4, which restates the theorem-like environments the retained text needs and adds the article's own macro definitions that the retained text needs (none), with extra wrapper packages (bm, bbm, dsfont, xspace, cancel, mathtools). The delivered numbering is the unit's own. Assets: no retained span contains a figure environment, a graphics-inclusion command, a PDF-page inclusion, a table or a TikZ picture; no image file is copied into this package, the package declares no asset dependency and no image file is redistributed with it. Licence: version arXiv:2411.10969v3 of this article is distributed under the Creative Commons Attribution 4.0 International licence, as recorded in the arXiv licence metadata for this exact version and shown on the abstract page https://arxiv.org/abs/2411.10969v3. A full rights notice is delivered at licenses/arxiv-2411.10969-CC-BY-4.0.txt. Characters: every retained excerpt is pure ASCII, so the delivered engine, XeLaTeX with its default Latin Modern text font, reports no missing character.
\begin{prop}\label{sec3:prop_vanishing}
If $f, g$ are smooth functions on $M$ such that $(\mathrm{d}f)_z,\,(\mathrm{d}g)_z\in \mathrm{ann}\bigl(\mathcal{D}_{\rho}(z)\bigr)$ at each $z\in M_{\mu}$, then 
$(\mathrm{d}\{f,\,g\}_{\Pi})_z\in \mathrm{ann}\bigl(\mathcal{D}_{\rho}(z)\bigr)$. 
\end{prop}

\begin{thm}\label{sec3:main thm}
Assume that the orbit space $\mathcal{M}_{\rho}$ is a smooth manifold such that $\pi_{\rho}$ is a submersion. 
If 
\begin{equation}\label{sec3:eqn_main}
\Pi^{\sharp}_z\bigl(\mathrm{ann}(\mathcal{D}_{\rho}(z))\bigr) \subset T_zM_{\mu} + \mathcal{D}_{\rho}(z)
\end{equation}
holds for each $z\in M_{\mu}$, then 
$\mathcal{M}_{\rho}$ is a Poisson manifold whose Poisson bracket $\{\cdot,\,\cdot\}_{\Pi/}$ is uniquely determined by 
\begin{equation}\label{sec3:eqn_main2}
 \{f,\,g\}_{\Pi/}\circ\pi_{\rho} = \{\tilde{f},\,\tilde{g}\}_{\Pi}\circ\iota 
\end{equation}
for any smooth function $f,g$ on $\mathcal{M}_{\rho}$, where $\tilde{f},\tilde{g}$ are extensions of $f\circ\pi_{\rho},\,g\circ\pi_{\rho}$, respectively, with conditions that 
$(\mathrm{d}\tilde{f})_z,\,(\mathrm{d}\tilde{g})_z\in \mathrm{ann}\bigl(\mathcal{D}_{\rho}(z)\bigr)$ at each $z\in M_{\mu}$. 

Conversely, if $\mathcal{M}_{\rho}$ is endowed with a Poisson structure by \eqref{sec3:eqn_main2}, then the condition \eqref{sec3:eqn_main} is satisfied. 
\end{thm}

\begin{proof}
Assume that \eqref{sec3:eqn_main} holds for any $z\in M_{\mu}$. Let $f,g$ be any smooth function on $\mathcal{M}_{\rho}$ 
and $\tilde{f}, \tilde{g}$ extensions of $f\circ\pi_{\rho}, g\circ\pi_{\rho}$. 
Namely, $\tilde{f}$ and $\tilde{g}$ are smooth functions on $M$ satisfying $\tilde{f}|_{M_{\mu}}=f\circ\pi_{\rho}$ and $\tilde{g}|_{M_{\mu}}=g\circ\pi_{\rho}$, respectively. 
Suppose that their differentials $\mathrm{d}\tilde{f},\,\mathrm{d}\tilde{g}$ vanish on $\mathcal{D}_{\rho}|_{M_{\mu}}$, and 
define a function $\{f,g\}_{\Pi/}$ on $\mathcal{M}_{\rho}$ as \eqref{sec3:eqn_main2}. 
Then, from Proposition \ref{sec3:prop_vanishing}, it follows that $\mathrm{d}\{\tilde{f},\,\tilde{g}\}_{\Pi}$ also vanishes on $\mathcal{D}_{\rho}|_{M_{\mu}}$. 
This implies that $\{\tilde{f},\,\tilde{g}\}_{\Pi}$ is constant on each leaf of $\mathcal{D}_{\rho}|_{M_{\mu}}$. Therefore, $\{f,g\}_{\Pi/}$ is well-defined as a function on $\mathcal{M}_{\rho}$. 
Next, we shall show that the function $\{f,g\}_{\Pi/}$ does not depend on the choice of their extensions. 
Let $\tilde{g}'$ be another extension of $g\circ\pi$ satisfying $(\mathrm{d}\tilde{g}')_z\in \mathrm{ann}\bigl(\mathcal{D}_{\rho}(z)\bigr)$ at each $z\in M_{\mu}$. 
Clearly, $\tilde{g}-\tilde{g}'=0$ on $M_{\mu}$, and $\mathrm{d}(\tilde{g}-\tilde{g}')$ vanishes on $\mathcal{D}_{\rho}|_{M_{\mu}}$. 
Therefore, using the assumption \eqref{sec3:eqn_main}, we have 
\[
\forall z\in M_{\mu};\ \bigl\langle (\mathrm{d}\tilde{f})_z,\, \Pi^{\sharp}_z\bigl(\mathrm{d}(\tilde{g}-\tilde{g}')_z\bigr)\bigr\rangle = 0, 
\]
alternatively 
\[
\forall z\in M_{\mu};\ \{\tilde{f},\,\tilde{g}\}_{\Pi}(z) = \{\tilde{f},\,\tilde{g}'\}_{\Pi}(z). 
\]
This means that the function $\{f,\,g\}_{\Pi/}\circ\pi_{\rho}$ is independent of the choice of extension. From the fact that $\{\cdot,\,\cdot\}_{\Pi}$ is a Poisson bracket, 
we see that $\mathcal{M}_{\rho}$ inherits a Poisson structure $\{\cdot,\,\cdot\}_{\Pi/}$ which satisfies $\pi_{\rho}^*\{\cdot,\,\cdot\}_{\Pi/}=\imath^*\{\cdot,\,\cdot\}_{\Pi}$. 

Conversely, we assume that the orbit space $\mathcal{M}_{\rho}$ is endowed with a Poisson structure determined by \eqref{sec3:eqn_main2}. 
We shall show the relation \eqref{sec3:eqn_main} or, equivalently, 
\begin{equation}\label{sec3:eqn_proof}
\mathrm{ann}\bigl(T_zM_{\mu}\bigr) \cap \mathrm{ann}\bigl(\mathcal{D}_{\rho}(z)\bigr)\subset \mathrm{ann}\Bigl(\Pi_z^{\sharp}\bigl(\mathrm{ann}\bigl(\mathcal{D}_{\rho}(z)\bigr)\bigr)\Bigr) 
\end{equation}
for any point $z$ in $M_{\mu}$. Fix any point $z$ in $M_{\mu}$. 
If $U_z$ is a slice chart for $M_{\mu}$ around $z$, then $\mathcal{D}_{\rho}$ is locally expressed as $\mathcal{D}_{\rho}=(U_z\cap M_{\mu})\times V$ with $V$ a vector subspace of 
$\mathbb{R}^k\oplus\mathbb{R}^{m-k}\cong \mathbb{R}^m$, where $m, k~(k<m)$ are dimensions of $M,\,M_{\mu}$, respectively. 

Take an element $\alpha_z\in \mathrm{ann}\bigl(T_zM_{\mu}\bigr) \cap \mathrm{ann}\bigl(\mathcal{D}_{\rho}(z)\bigr)$ to be arbitrary. 
From $\alpha_z\in \mathrm{ann}\bigl(T_zM_{\mu}\bigr)$, 
$\alpha_z$ can be written locally as $\alpha_z = (\mathrm{d}F)_z$  %= a_1(\mathrm{d}x_{k+1})_z + a_2(\mathrm{d}x_{k+2})_z+\cdots +a_{m-k}(\mathrm{d}x_m)_z 
with $F\in C^{\infty}(U_z)$ defined as 
\[
F(\boldsymbol{x})=a_1x_{k+1} + a_2x_{k+2}+\cdots +a_{m-k}x_{m},\quad \boldsymbol{x}\in U_z
\]
for some $\boldsymbol{a}=(a_1,\cdots,a_{m-k})\in \mathbb{R}^{m-k}$. Moreover, since $\alpha_z\in \mathrm{ann}\bigl(\mathcal{D}_{\rho}(z)\bigr)$, it turns out that $F$ satisfies 
\begin{equation}\label{sec3:eqn_proof2}
 F|_{U_z\cap M_{\mu}} = 0\quad \text{and}\quad \forall y\in U_z\cap M_{\mu};\ (\mathrm{d}F)_y|_{\mathcal{D}_{\rho}(y)} = \boldsymbol{0}. 
\end{equation}

Next, let $\beta_z\in \mathrm{ann}\bigl(\mathcal{D}_{\rho}(z)\bigr)$ and $G$ a function on $U_z$ such that 
\[
\beta_z=(\mathrm{d}G)_z\quad \text{and}\quad \forall y\in U_z\cap M_{\mu};\ 
(\mathrm{d}G)_y|_{\mathcal{D}_{\rho}(y)} = \boldsymbol{0}. 
\]
If $g$ is a function on $\mathcal{M}_{\rho}$ such that the extension of which is $G$, then 
\[
 \langle \alpha_z,\,\Pi^{\sharp}_z(\beta_z)\rangle = \{F,\,G\}_{\Pi}(z)=\{0,\,g\}_{\Pi/}\circ \pi_{\rho}(z)=0. 
\]
This means that $\alpha_z\in \mathrm{ann}\bigl(\Pi_z^{\sharp}\bigl(\mathrm{ann}\bigl(\mathcal{D}_{\rho}(z)\bigr)\bigr)\bigr)$. Therefore, \eqref{sec3:eqn_proof} holds. 
\end{proof}

\end{document}
